% GENERATED FILE. Edit canonical lessons, then run npm run generate:books.
% canonical-source-hash: fnv1a64:f3ed0ba730499e8a
% canonical-lessons: MR-C459-odh, MR-C459-hurhur, MR-C459-aturta, MR-C459-asusne, MR-C459-jhurne

\chapter{Longing}
\label{ch:mr-longing}
\hlchaptermodality{459}

\begin{chapteropening}
\textbf{By the end of this chapter:} \emph{I can talk about longing, eagerness and pining.}

Complete the last lesson of chapter 459: I can talk about longing, eagerness and pining.
\end{chapteropening}

\section[\texorpdfstring{\mr{ओढ} (oḍh)}{ओढ (oḍh)}]{\mr{ओढ} (oḍh) — a longing; a pull}

\label{lesson:MR-C459-odh}



\begin{quote}
Before the new one: say the Marathi for to feel relieved, then the Marathi for to be choked up.
\end{quote}

\subsection*{You'll want to know: \mr{ओढ}}
\textbf{\mr{ओढ}} — \emph{oḍh} — \textquotedblleft{}a longing; a pull\textquotedblright{}.

\begin{quote}
\mr{सुट्टीत} \mr{गावाची} \mr{ओढ} \mr{लागते}.
\end{quote}

\begin{quote}
\emph{suṭṭīt gāvācī oḍh lāgte.}
\end{quote}

\begin{quote}
In the holidays one longs for the village.
\end{quote}

Feminine.

\mr{ओढ} \mr{लागणे}, to feel drawn to, with the genitive: \mr{घराची} \mr{ओढ}. It belongs with \mr{ओढणे}, to pull.

\begin{cousinweb}
\textbf{\mr{ओढ}} (\emph{oḍh}) is a Marathi noun beside \mr{ओढणे}, \textquotedblleft{}to pull\textquotedblright{}, which goes back to the reconstructed \emph{*avakaḍḍhati}, \textquotedblleft{}pulls away\textquotedblright{} (CDIAL 14245).
\end{cousinweb}

\begin{grammarlens}[title={What you've built}]
One more word for describing what you lived and felt.
\end{grammarlens}

\subsection*{Your turn}
\begin{itemize}
\raggedright
  \item \emph{Say it:} \emph{oḍh}
  \item \emph{Say it:} \emph{oḍh}, once more
  \item \emph{Say it:} \emph{suṭṭīt gāvācī oḍh lāgte.} — In the holidays one longs for the village.
  \item \emph{Recall:} say the Marathi for to feel relieved, then the Marathi for to be choked up, then say \emph{oḍh} again
\end{itemize}

\begin{tcolorbox}[breakable,colback=teal!4,colframe=teal!35!black,title={Before you move on}]
What does \mr{ओढ} mean? (\textquotedblleft{}A longing; a pull\textquotedblright{}.) Say it once more.
\end{tcolorbox}
\section[\texorpdfstring{\mr{हुरहूर} (hurhūr)}{हुरहूर (hurhūr)}]{\mr{हुरहूर} (hurhūr) — a wistful unease}

\label{lesson:MR-C459-hurhur}



\begin{quote}
Before the new one: say the Marathi for to be choked up, then the Marathi for a longing.
\end{quote}

\subsection*{You'll want to know: \mr{हुरहूर}}
\textbf{\mr{हुरहूर}} — \emph{hurhūr} — \textquotedblleft{}a wistful unease\textquotedblright{}.

\begin{quote}
\mr{शाळेचा} \mr{शेवटचा} \mr{दिवस} \mr{होता}, \mr{मनात} \mr{हुरहूर} \mr{होती}.
\end{quote}

\begin{quote}
\emph{śāḷecā śevaṭcā divas hotā, manāt hurhūr hotī.}
\end{quote}

\begin{quote}
It was the last day of school, and there was an ache in the heart.
\end{quote}

Feminine.

The ache at an ending or a parting: \mr{हुरहूर} \mr{लागणे}, with the dative.

\begin{cousinweb}
\textbf{\mr{हुरहूर}} (\emph{hurhūr}) is a Marathi noun built on a doubled syllable; its origin is not settled.
\end{cousinweb}

\begin{grammarlens}[title={What you've built}]
One more word for describing what you lived and felt.
\end{grammarlens}

\subsection*{Your turn}
\begin{itemize}
\raggedright
  \item \emph{Say it:} \emph{hurhūr}
  \item \emph{Say it:} \emph{hurhūr}, once more
  \item \emph{Say it:} \emph{śāḷecā śevaṭcā divas hotā, manāt hurhūr hotī.} — It was the last day of school, and there was an ache in the heart.
  \item \emph{Recall:} say the Marathi for to be choked up, then the Marathi for a longing, then say \emph{hurhūr} again
\end{itemize}

\begin{tcolorbox}[breakable,colback=teal!4,colframe=teal!35!black,title={Before you move on}]
What does \mr{हुरहूर} mean? (\textquotedblleft{}A wistful unease\textquotedblright{}.) Say it once more.
\end{tcolorbox}
\section[\texorpdfstring{\mr{आतुरता} (āturtā)}{आतुरता (āturtā)}]{\mr{आतुरता} (āturtā) — eagerness}

\label{lesson:MR-C459-aturta}



\begin{quote}
Before the new one: say the Marathi for a longing, then the Marathi for a wistful unease.
\end{quote}

\subsection*{You'll want to know: \mr{आतुरता}}
\textbf{\mr{आतुरता}} — \emph{āturtā} — \textquotedblleft{}eagerness\textquotedblright{}.

\begin{quote}
\mr{निकालाची} \mr{सर्वांना} \mr{आतुरता} \mr{होती}.
\end{quote}

\begin{quote}
\emph{nikālācī sarvānnā āturtā hotī.}
\end{quote}

\begin{quote}
Everyone was eager for the result.
\end{quote}

Feminine.

\mr{आतुर}, eager, + -\mr{ता}; \mr{आतुरतेने} \mr{वाट} \mr{पाहणे}, to wait eagerly.

\begin{cousinweb}
\textbf{\mr{आतुरता}} (\emph{āturtā}) is built on the tatsama \mr{आतुर}, Sanskrit \emph{ātura}, glossed \textquotedblleft{}suffering, sick\textquotedblright{} (Monier-Williams); Marathi \mr{आतुर} has moved to \textquotedblleft{}eager\textquotedblright{}, and \emph{-tā} makes the noun.
\end{cousinweb}

\begin{grammarlens}[title={What you've built}]
One more word for describing what you lived and felt.
\end{grammarlens}

\subsection*{Your turn}
\begin{itemize}
\raggedright
  \item \emph{Say it:} \emph{āturtā}
  \item \emph{Say it:} \emph{āturtā}, once more
  \item \emph{Say it:} \emph{nikālācī sarvānnā āturtā hotī.} — Everyone was eager for the result.
  \item \emph{Recall:} say the Marathi for a longing, then the Marathi for a wistful unease, then say \emph{āturtā} again
\end{itemize}

\begin{tcolorbox}[breakable,colback=teal!4,colframe=teal!35!black,title={Before you move on}]
What does \mr{आतुरता} mean? (\textquotedblleft{}Eagerness\textquotedblright{}.) Say it once more.
\end{tcolorbox}
\section[\texorpdfstring{\mr{आसुसणे} (āsusṇe)}{आसुसणे (āsusṇe)}]{\mr{आसुसणे} (āsusṇe) — to long for}

\label{lesson:MR-C459-asusne}



\begin{quote}
Before the new one: say the Marathi for a wistful unease, then the Marathi for eagerness.
\end{quote}

\subsection*{You'll want to know: \mr{आसुसणे}}
\textbf{\mr{आसुसणे}} — \emph{āsusṇe} — \textquotedblleft{}to long for\textquotedblright{}.

\begin{quote}
\mr{पावसासाठी} \mr{सगळं} \mr{गाव} \mr{आसुसलं} \mr{होतं}.
\end{quote}

\begin{quote}
\emph{pāvsāsāṭhī sagḷaṁ gāv āsuslaṁ hotaṁ.}
\end{quote}

\begin{quote}
The whole village was longing for rain.
\end{quote}

Past (he …): \mr{तो} \mr{आसुसला}, \emph{to āsuslā}.

Takes -\mr{साठी}, or the -\mr{ायला} form: \mr{भेटायला} \mr{आसुसणे}, to long to meet.

\begin{cousinweb}
\textbf{\mr{आसुसणे}} (\emph{āsusṇe}) is a Marathi verb whose origin is not settled.
\end{cousinweb}

\begin{grammarlens}[title={What you've built}]
One more word for describing what you lived and felt.
\end{grammarlens}

\subsection*{Your turn}
\begin{itemize}
\raggedright
  \item \emph{Say it:} \emph{āsusṇe}
  \item \emph{Say it:} \emph{āsusṇe}, once more
  \item \emph{Say it:} \emph{pāvsāsāṭhī sagḷaṁ gāv āsuslaṁ hotaṁ.} — The whole village was longing for rain.
  \item \emph{Recall:} say the Marathi for a wistful unease, then the Marathi for eagerness, then say \emph{āsusṇe} again
\end{itemize}

\begin{tcolorbox}[breakable,colback=teal!4,colframe=teal!35!black,title={Before you move on}]
What does \mr{आसुसणे} mean? (\textquotedblleft{}To long for\textquotedblright{}.) Say it once more.
\end{tcolorbox}
\section[\texorpdfstring{\mr{झुरणे} (jhurṇe)}{झुरणे (jhurṇe)}]{\mr{झुरणे} (jhurṇe) — to pine}

\label{lesson:MR-C459-jhurne}



\begin{quote}
Before the new one: say the Marathi for eagerness, then the Marathi for to long for.
\end{quote}

\subsection*{You'll want to know: \mr{झुरणे}}
\textbf{\mr{झुरणे}} — \emph{jhurṇe} — \textquotedblleft{}to pine\textquotedblright{}.

\begin{quote}
\mr{दूर} \mr{गेलेल्या} \mr{मुलासाठी} \mr{आई} \mr{झुरत} \mr{होती}.
\end{quote}

\begin{quote}
\emph{dūr gelelyā mulāsāṭhī āī jhurat hotī.}
\end{quote}

\begin{quote}
The mother was pining for her son who had gone far away.
\end{quote}

Past (he …): \mr{तो} \mr{झुरला}, \emph{to jhurlā}.

Slow and wearing: \mr{एखाद्यासाठी} \mr{झुरणे}, to pine for someone.

\begin{cousinweb}
\textbf{\mr{झुरणे}} (\emph{jhurṇe}) is inherited: CDIAL lists Marathi \emph{jhurṇẽ}, \textquotedblleft{}to waste away\textquotedblright{}, under the reconstructed \emph{*jhurati} (CDIAL 5405).
\end{cousinweb}

\begin{grammarlens}[title={What you've built}]
One more word for describing what you lived and felt.
\end{grammarlens}

\subsection*{Your turn}
\begin{itemize}
\raggedright
  \item \emph{Say it:} \emph{jhurṇe}
  \item \emph{Say it:} \emph{jhurṇe}, once more
  \item \emph{Say it:} \emph{dūr gelelyā mulāsāṭhī āī jhurat hotī.} — The mother was pining for her son who had gone far away.
  \item \emph{Recall:} say the Marathi for eagerness, then the Marathi for to long for, then say \emph{jhurṇe} again
\end{itemize}

\begin{tcolorbox}[breakable,colback=teal!4,colframe=teal!35!black,title={Before you move on}]
What does \mr{झुरणे} mean? (\textquotedblleft{}To pine\textquotedblright{}.) Say it once more.
\end{tcolorbox}
